%----------------------------------------------------------------------------------------
% NOTATIE EN FIGUREN VAN DE CURSUS EIGENWAARDEN EN EIGENVECTOREN
% Staat in de preambule en niet in een hoofdstuk: ook de lijsten met oplossingen gebruiken dit.
%----------------------------------------------------------------------------------------
% Bewaarde grafiek, geplaatst door maak-pdf.py: bestand en pagina in figuren-cache/
\newcommand{\grafiekbestand}[2]{\includegraphics[page=#2,trim=3bp 3bp 3bp 3bp]{figuren-cache/#1}}

\providecommand{\N}{\mathbb{N}}
\providecommand{\Z}{\mathbb{Z}}
\providecommand{\Q}{\mathbb{Q}}
\providecommand{\R}{\mathbb{R}}
\providecommand{\C}{\mathbb{C}}
\newcommand{\e}{\mathrm{e}}

% Zoals stappen, maar met een bredere kolom voor de berekening: voor lange formules
\newenvironment{stappenbreed}{%
  \par\renewcommand{\arraystretch}{1.25}%
  \noindent\begin{tabular}{@{}>{$\displaystyle}p{9.8cm}<{$}@{\hspace{4mm}}>{\RaggedRight\small\itshape\arraybackslash}p{4.9cm}@{}}%
}{\end{tabular}\par}

% Zoals stappen, met eigen kolombreedtes: \begin{stappenvrij}{6.2cm}{3.8cm}. Voor naast een figuur.
\newenvironment{stappenvrij}[2]{%
  \par\renewcommand{\arraystretch}{1.25}%
  \noindent\begin{tabular}{@{}>{$\displaystyle}p{#1}<{$}@{\hspace{3mm}}>{\RaggedRight\small\itshape\arraybackslash}p{#2}@{}}%
}{\end{tabular}\par}

% Tekst of berekening links, figuur rechts: \naast[breedte rechts]{links}{rechts}
% De breedte rechts komt uit de figuur zelf; het optionele argument wordt niet gebruikt.
\newsavebox{\naastbox}
\newcommand{\naast}[3][]{%
  \par\noindent\sbox{\naastbox}{#3}%
  \begin{minipage}[t]{\dimexpr\linewidth-\wd\naastbox-3mm\relax}\vspace{0pt}#2\end{minipage}\hfill
  \begin{minipage}[t]{\wd\naastbox}\vspace{0pt}\usebox{\naastbox}\end{minipage}\par\addvspace{1mm}}
% Een rij figuren, gecentreerd
\newcommand{\figuren}[1]{\par\addvspace{1mm}{\centering #1\par}\addvspace{1.5mm}}
% Opgave met een grafiek rechts: \begin{metgrafiek}[breedte]{tikz} opgave en deelvragen \end{metgrafiek}
\newsavebox{\grafiekbox}
\newenvironment{metgrafiek}[2][]{%
  \sbox{\grafiekbox}{#2}%
  \begin{minipage}[t]{\dimexpr\linewidth-\wd\grafiekbox-3mm\relax}%
}{%
  \end{minipage}\hfill
  \begin{minipage}[t]{\wd\grafiekbox}\vspace{-2mm}\usebox{\grafiekbox}\end{minipage}\par}

% Rooster met assen: \rooster{xmin}{xstap}{xmax}{ymin}{ystap}{ymax}, met getallen op de assen (gehele waarden).
% \roosterleeg doet hetzelfde zonder getallen.
\newcommand{\roosterleeg}[6]{%
  \draw[black!14,very thin,xstep=#2,ystep=#5] (#1,#4) grid (#3,#6);
  \draw[->,grijs] (#1,0) -- ({#3+.4*#2},0) node[right,text=black] {$x$};
  \draw[->,grijs] (0,#4) -- (0,{#6+.4*#5}) node[above,text=black] {$y$};}
\newcommand{\rooster}[6]{%
  \roosterleeg{#1}{#2}{#3}{#4}{#5}{#6}
  \pgfmathtruncatemacro{\rtweex}{#1+#2}\pgfmathtruncatemacro{\rtweey}{#4+#5}
  \foreach \t in {#1,\rtweex,...,#3}{\ifnum\t=0\else\node[below,inner sep=1.5pt,text=black] at (\t,0) {$\t$};\fi}
  \foreach \t in {#4,\rtweey,...,#6}{\ifnum\t=0\else\node[left,inner sep=1.5pt,text=black] at (0,\t) {$\t$};\fi}}
% Assen zonder rooster: \assen{xmin}{xmax}{ymin}{ymax}, en een kromme: \kromme[opties]{van}{tot}{voorschrift in \x}
\newcommand{\assen}[4]{%
  \draw[->,lijngrijs] (#1,0) -- (#2,0) node[right,text=black] {$x$};
  \draw[->,lijngrijs] (0,#3) -- (0,#4) node[above,text=black] {$y$};}
\newcommand{\kromme}[4][]{\draw[accent,thick,#1] plot[domain=#2:#3,samples=90,smooth] (\x,{#4});}

% Ruimtelijk assenstelsel in schuine projectie: gebruik [ruimte] als optie van tikzpicture,
% en \assendrie{lengte}. De x-as wijst schuin naar voren, y naar rechts, z naar boven.
\tikzset{ruimte/.style={x={(-0.45cm,-0.38cm)},y={(1cm,0cm)},z={(0cm,1cm)}}}
\newcommand{\assendrie}[1]{%
  \draw[->,lijngrijs] (0,0,0) -- (#1,0,0) node[below left,text=black] {$x$};
  \draw[->,lijngrijs] (0,0,0) -- (0,#1,0) node[right,text=black] {$y$};
  \draw[->,lijngrijs] (0,0,0) -- (0,0,#1) node[above,text=black] {$z$};}

% Uit de cursus Matrices: compacte matrix \km{1 & 2 \\ 3 & 4}, uitgebreide matrix \um{cc|c}{...},
% en een rijbewerking naast een matrix \rb{R_2 - 2R_1}.
\newcommand{\km}[1]{{\setlength{\arraycolsep}{2.4pt}\renewcommand{\arraystretch}{.82}\left[\begin{array}{*{10}{c}}#1\end{array}\right]}}
\newcommand{\um}[2]{{\setlength{\arraycolsep}{4pt}\left[\begin{array}{#1}#2\end{array}\right]}}
\newcommand{\rb}[1]{\text{\small\color{grijs}$#1$}}

%----------------------------------------------------------------------------------------
% EIGENWAARDEN EN EIGENVECTOREN (notatie zoals in Vectorruimten en Multivariate statistiek)
%----------------------------------------------------------------------------------------
% Kolomvector: \kol{1 \\ 2 \\ 3}
\newcommand{\kol}[1]{{\setlength{\arraycolsep}{3pt}\begin{bmatrix} #1 \end{bmatrix}}}
\newcommand{\rang}{\operatorname{rang}}
\newcommand{\Ker}{\operatorname{Ker}}
\newcommand{\nulv}{\vec{0}}
% De deelruimte voortgebracht door vectoren: \vb{\vec u, \vec v}
\newcommand{\vb}[1]{\left\langle #1 \right\rangle}
% Getransponeerde: A\tr
\newcommand{\tr}{^{\mkern-1.5mu\mathsf{T}}}
% Spoor en diagonaalmatrix: \spoor A, \diag(2, 3)
\newcommand{\spoor}{\operatorname{sp}}
\newcommand{\diag}{\operatorname{diag}}
% Eigenruimte bij lambda: \Eig{\lambda}
\newcommand{\Eig}[1]{E_{#1}}

% FIGUREN, binnen een tikzpicture
% Een vector als pijl vanuit de oorsprong of een ander punt: \vpijl[opties]{(van)}{(naar)}
\newcommand{\vpijl}[3][]{\draw[accent,very thick,-{Stealth[length=2.2mm]},#1] #2 -- #3;}
% Het beeld van een vector, dunner en fel: \bpijl[opties]{(van)}{(naar)}
\newcommand{\bpijl}[3][]{\draw[fel!85!black,thick,-{Stealth[length=2mm]},#1] #2 -- #3;}
% Een rechte door de oorsprong in de richting (a, b), gestippeld, van t = -s tot t = s: \richting{a}{b}{s}
\newcommand{\richting}[3]{\draw[accent!55,dashed] ({-(#3)*(#1)},{-(#3)*(#2)}) -- ({(#3)*(#1)},{(#3)*(#2)});}
% Scheef rooster van een basis u, v in het vlak: \scheefrooster{u1}{u2}{v1}{v2}{imin}{imax}{jmin}{jmax}
\newcommand{\scheefrooster}[8]{%
  \foreach \i in {#5,...,#6}{\draw[accent!30,thin] ({\i*(#1)+(#7)*(#3)},{\i*(#2)+(#7)*(#4)}) -- ({\i*(#1)+(#8)*(#3)},{\i*(#2)+(#8)*(#4)});}
  \foreach \j in {#7,...,#8}{\draw[accent!30,thin] ({(#5)*(#1)+\j*(#3)},{(#5)*(#2)+\j*(#4)}) -- ({(#6)*(#1)+\j*(#3)},{(#6)*(#2)+\j*(#4)});}}
% Grafen uit Matrices (hoofdstuk 4): \node[knoop] (a) at (0,0) {A}; \draw[boog] (a) -- (b); kans bij een boog: node[kans] {0,3}
\tikzset{
  knoop/.style={circle,draw=accent,fill=accent!12,thick,minimum size=7mm,inner sep=0pt,font=\small\bfseries},
  boog/.style={-{Stealth[length=2.2mm]},accent,thick,shorten >=1pt,shorten <=1pt},
  lus/.style={out=#1+25,in=#1-25,looseness=7,min distance=7mm},
  lus/.default=90,
  kans/.style={fill=white,inner sep=1pt,font=\scriptsize,text=black},
  hulplijn/.style={grijs,thick,densely dotted},
}
% Absolute waarde: \abs{x}
\providecommand{\abs}[1]{\lvert #1\rvert}
